\documentclass[12pt,a4paper]{article}

\usepackage{fontspec}
\setmainfont{Times New Roman}
\setmonofont{Menlo}
\usepackage[russian]{babel}

\setlength{\oddsidemargin}{4.6mm}
\setlength{\evensidemargin}{4.6mm}
\setlength{\textwidth}{165mm}
\setlength{\topmargin}{-5.4mm}
\setlength{\headheight}{0pt}
\setlength{\headsep}{0pt}
\setlength{\textheight}{250mm}
\setlength{\footskip}{12mm}
\setlength{\parindent}{12mm}
\linespread{1.2}
\sloppy

\begin{document}
\special{papersize=210mm,297mm}

\begin{titlepage}
\begin{center}
\textbf{Московский авиационный институт\\
(национальный исследовательский университет)}

\vspace{55mm}

\textbf{ОТЧЁТ\\
по лабораторной работе № 5\\
по дисциплине «Операционные системы»}

\vspace{8mm}

Тема: «Динамические библиотеки»\\
Вариант 1
\end{center}

\vspace{45mm}

\begin{flushright}
\begin{tabular}{ll}
Выполнил: & Абазатов Роман Аленович \\
Группа: & М8О-202БВ-25 \\
Проверил: & преподаватель Сысоев М.А. \\
Дата: & 23.11.2026 \\
\end{tabular}
\end{flushright}

\vfill

\begin{center}
Москва, 2026
\end{center}
\end{titlepage}

\section{Цель работы}

Получить практические навыки создания динамических библиотек и программ,
которые используют их функции. Сравнить подключение библиотеки на этапе
компоновки и загрузку библиотеки во время выполнения программы.

\section{Постановка задачи}

Для варианта 1 требуется реализовать две функции с одинаковыми именами и
сигнатурами в двух динамических библиотеках:

\begin{verbatim}
float SinIntegral(float A, float B, float e);
float Derivative(float A, float delta_x);
\end{verbatim}

Функция \texttt{SinIntegral} приближённо вычисляет интеграл \texttt{sin(x)}
от \texttt{A} до \texttt{B} с шагом \texttt{e}. В первой библиотеке используется
метод левых прямоугольников, во второй --- метод трапеций.

Функция \texttt{Derivative} приближённо вычисляет производную \texttt{cos(x)}
в точке \texttt{A} при приращении \texttt{delta\_x}. Первая библиотека
использует правую конечную разность:

\begin{verbatim}
(cos(A + delta_x) - cos(A)) / delta_x
\end{verbatim}

Вторая использует центральную конечную разность:

\begin{verbatim}
(cos(A + delta_x) - cos(A - delta_x)) / (2 * delta_x)
\end{verbatim}

Программа № 1 должна вызывать функции библиотеки, выбранной при компоновке.
Программа № 2 должна загружать библиотеки по относительным путям во время
выполнения. Команда \texttt{0} переключает реализации в программе № 2;
команда \texttt{1 A B e} вызывает первую функцию, а команда
\texttt{2 A deltaX} --- вторую.

\section{Схема решения}

Обе библиотеки экспортируют один контракт, объявленный в
\texttt{include/contract.h}. Различаются только численные алгоритмы.
Схема подключения библиотек:

\begin{center}
\begin{tabular}{ccl}
\texttt{program1} & $\longrightarrow$ & \texttt{libimpl1.so} (компоновка) \\
\texttt{program2} & $\longrightarrow$ & \texttt{libimpl1.so} (\texttt{dlopen}) \\
команда \texttt{0} & $\longrightarrow$ & \texttt{libimpl2.so} (\texttt{dlopen}) \\
\end{tabular}
\end{center}

Программа № 1 связана с \texttt{libimpl1.so}; загрузчик Linux открывает
библиотеку при запуске. В её исполняемом файле установлен
\texttt{RUNPATH=\$ORIGIN}, поэтому библиотека ищется рядом с программой.

Программа № 2 сначала открывает \texttt{./build/libimpl1.so} через
\texttt{dlopen}. Адреса функций она получает через \texttt{dlsym}. При
переключении сначала загружается другая библиотека и проверяются обе функции,
затем закрывается прежняя. Если загрузка не удалась, текущая реализация
остаётся активной. Относительные пути требуют запуска программы из каталога
\texttt{lab\_5}.

\section{Реализация}

\subsection{Интегрирование и дифференцирование}

При интегрировании интервал делится на участки длиной не более \texttt{e}.
Последний участок может быть короче. Для левого прямоугольника площадь
участка равна значению \texttt{sin(x)}, умноженному на его ширину. Для трапеции
берётся среднее значений функции на двух концах участка:

\begin{verbatim}
for (double x = left; x < right;) {
    double next = fmin(x + (double)e, right);
    if (next <= x) {
        errno = ERANGE;
        return NAN;
    }
    sum += (sin(x) + sin(next)) * (next - x) / 2.0;
    x = next;
}
\end{verbatim}

Это фрагмент реализации метода трапеций. Для прямоугольников внутри цикла
используется выражение \texttt{sin(x) * (next - x)}. Если \texttt{A > B}, знак
результата меняется. Внутри библиотеки вычисления выполняются с типом
\texttt{double}, затем результат приводится к \texttt{float} согласно
контракту. Неконечные аргументы, неположительный шаг и нулевое приращение
производной отклоняются.

\subsection{Подключение библиотек}

В программе № 1 функции объявлены в заголовочном файле и вызываются
непосредственно. При сборке указывается флаг \texttt{-limpl1}.

Программа № 2 загружает выбранную библиотеку и находит два символа:

\begin{verbatim}
void *handle = dlopen(path, RTLD_NOW | RTLD_LOCAL);
sin_integral_fn integral =
    (sin_integral_fn)dlsym(handle, "SinIntegral");
derivative_fn derivative =
    (derivative_fn)dlsym(handle, "Derivative");
\end{verbatim}

После каждого вызова \texttt{dlsym} проверяется \texttt{dlerror}. Адреса
функций используются только после успешной загрузки обоих символов. Функция
\texttt{dlclose} освобождает прежний объект после переключения и текущий
объект перед завершением программы.

Общий для двух программ модуль \texttt{src/cli.c} считывает команды построчно,
проверяет число аргументов и их значения. Некорректная строка не завершает
программу; сообщение выводится в \texttt{stderr}.

\section{Сборка и проверка}

Сборка и проверка выполнены 23 ноября 2026 года в Linux-контейнере
на архитектуре aarch64: GCC 14.4.0, glibc 2.41. Для воспроизведения:

\begin{verbatim}
cd lab_5
make clean all
printf '1 0 1 0.25\n2 1 0.1\n' | ./build/program1
printf '1 0 1 0.25\n0\n1 0 1 0.25\n' | ./build/program2
\end{verbatim}

Для сравнения методов интеграл вычислялся на отрезке от 0 до приближённого
значения числа π с шагом 0,001. Его точное значение близко к 2. Производная
\texttt{cos(x)} в точке 1 равна \texttt{-sin(1)}, то есть приблизительно
\texttt{-0.841470985}.

\begin{center}
\small
\begin{tabular}{|p{62mm}|r|r|}
\hline
Команда & Реализация 1 & Реализация 2 \\
\hline
\texttt{1 0 3.1415927 0.001} & 1.99999976 & 1.99999988 \\
\hline
\texttt{2 1 0.001} & -0.841741025 & -0.841470838 \\
\hline
\end{tabular}
\end{center}

Пример работы переключения реализаций:

\begin{verbatim}
$ printf '1 0 1 0.25\n0\n1 0 1 0.25\n0\n1 0 1 0.25\n' \
    | ./build/program2
0.352117062
Реализация 2
0.457300931
Реализация 1
0.352117062
\end{verbatim}

Дополнительно проверены обратный порядок границ интегрирования, нулевой шаг,
нулевое приращение, пропущенные и лишние аргументы, а также текст вместо
числа. При отсутствии второй библиотеки \texttt{dlopen} выдаёт ошибку, после
которой первая библиотека остаётся доступной.

Команда \texttt{readelf -d} показала следующие зависимости:

\begin{verbatim}
program1: NEEDED libimpl1.so, libm.so.6, libc.so.6
          RUNPATH [$ORIGIN]
program2: NEEDED libm.so.6, libc.so.6
\end{verbatim}

В программе № 2 нет записи \texttt{NEEDED} для созданных библиотек: их имена
передаются \texttt{dlopen} во время исполнения. В данной версии glibc функции
\texttt{dlopen} входят в \texttt{libc}, поэтому отдельная запись о
\texttt{libdl.so} не требуется.

\section{Сравнение способов использования библиотек}

\begin{center}
\small
\begin{tabular}{|p{41mm}|p{50mm}|p{50mm}|}
\hline
Свойство & Программа № 1 & Программа № 2 \\
\hline
Выбор реализации & При компоновке & Во время работы по команде \texttt{0} \\
\hline
Вызов функции & Прямой вызов по контракту & Через указатель, полученный \texttt{dlsym} \\
\hline
Отсутствие библиотеки & Программа не запускается & Ошибка \texttt{dlopen} обрабатывается программой \\
\hline
Замена реализации & Требует другой сборки или замены файла библиотеки & Выполняется без перезапуска \\
\hline
\end{tabular}
\end{center}

Оба способа используют динамические библиотеки. При компоновке зависимость
записывается в исполняемый файл и обрабатывается загрузчиком Linux. При
явной загрузке программа сама выбирает файл библиотеки, получает адреса
функций и управляет временем её использования.

\section{Вывод}

Созданы две динамические библиотеки с единым контрактом и разными методами
вычисления интеграла и производной. Первая программа использует библиотеку,
выбранную при компоновке; вторая загружает библиотеки во время работы и
переключает реализации. Корректность проверена на обычных и ошибочных
командах в Linux. Различие способов подключения подтверждено содержимым
таблицы динамических зависимостей исполняемых файлов.

\end{document}
